% --- Schrift: Times-Klon (Nimbus Roman) ---
\usepackage{mathptmx}

% --- Kolumnentitel: Kapitelname in Kapitälchen links, Seitenzahl rechts ---
\usepackage{fancyhdr}

\pagestyle{fancy}
\fancyhf{}
\fancyhead[L]{\footnotesize\scshape\nouppercase{\rightmark}}
\fancyhead[R]{\footnotesize\thepage}
\renewcommand{\headrulewidth}{0pt}
\fancypagestyle{plain}{\fancyhf{}\renewcommand{\headrulewidth}{0pt}}

% Akt-Eröffnungsseite: Kopfzeilentitel unterdrücken (er duplizierte die
% große Überschrift direkt darunter), Seitenzahl bleibt.
\usepackage{etoolbox}
\fancypagestyle{aktseite}{%
  \fancyhf{}%
  \fancyhead[R]{\footnotesize\thepage}%
  \renewcommand{\headrulewidth}{0pt}%
}

% --- Artefakt-Rahmen: linke Linie ---
\usepackage[framemethod=TikZ]{mdframed}
\usepackage{needspace}

\newmdenv[
  topline=false, bottomline=false, rightline=false,
  leftline=true, linewidth=1pt, linecolor=black,
  innerleftmargin=12pt, innerrightmargin=0pt,
  innertopmargin=10pt, innerbottommargin=10pt,
  skipabove=\baselineskip, skipbelow=\baselineskip,
  nobreak=true
]{artefaktframe}

% artefaktkeep: darf NICHT über einen Seitenumbruch laufen
\newmdenv[
  topline=false, bottomline=false, rightline=false,
  leftline=true, linewidth=1pt, linecolor=black,
  innerleftmargin=12pt, innerrightmargin=0pt,
  innertopmargin=10pt, innerbottommargin=10pt,
  skipabove=\baselineskip, skipbelow=\baselineskip,
  nobreak=true
]{artefaktkeepframe}

% artefaktbreak: darf über einen Seitenumbruch laufen.
% Nur für das Abstract ‹7337› (Kap. 19), das sechs Absätze hat und damit
% länger ist als eine Seite; alle übrigen Artefakte bleiben ungeteilt.
\newmdenv[
  topline=false, bottomline=false, rightline=false,
  leftline=true, linewidth=1pt, linecolor=black,
  innerleftmargin=12pt, innerrightmargin=0pt,
  innertopmargin=10pt, innerbottommargin=10pt,
  skipabove=\baselineskip, skipbelow=\baselineskip,
  nobreak=false
]{artefaktbreakframe}

% Artefakte kursiv setzen. Kein \needspace hier: in Kombination mit
% mdframed (nobreak=true) erzeugte es eine zusätzliche Leerseite vor
% Artefakten, die direkt nach einem \clearpage beginnen.
% Vermerk „Fiktives Artefakt“: um 90° gedreht, außen links am Strich,
% unten bündig mit dem Strichende. Nimmt keinen eigenen Platz ein.
% Wird vom Filter (Pass 1) an jedes Artefakt angehängt.
\newcommand{\artefaktvermerk}{%
  \par\vspace{-\parskip}\nointerlineskip\noindent
  \llap{\raisebox{-11pt}[0pt][0pt]{\rotatebox{90}{%
    \upshape\scriptsize\scshape Fiktives Artefakt}}\hspace{17pt}}\par}

\newenvironment{artefakt}
  {\begin{artefaktframe}\itshape}
  {\end{artefaktframe}}
\newenvironment{artefaktkeep}
  {\begin{artefaktkeepframe}\itshape}
  {\end{artefaktkeepframe}}
\newenvironment{artefaktbreak}
  {\begin{artefaktbreakframe}\itshape}
  {\end{artefaktbreakframe}}

% --- Terminalabsätze: Courier, ohne Kursiv, ohne Silbentrennung ---
% Der Filter (Pass 0) setzt alles, was eine Maschine ausgibt, in diese Umgebung.
% #1 Einzug links, #2 Absatzabstand im Block, #3 Schriftgröße in pt.
\usepackage{courier}
\newlength{\termskipsave}
\newenvironment{term}[3]
  {\par\setlength{\termskipsave}{\parskip}\begingroup
   \ttfamily\fontsize{#3}{13.6}\selectfont
   \raggedright\leftskip=#1\relax\parindent=0pt\relax
   \hyphenpenalty=10000\exhyphenpenalty=10000\interlinepenalty=10000\relax
   \setlength{\parskip}{#2}\relax
   \vspace{\dimexpr\termskipsave-#2\relax}}
  {\par\endgroup}

% --- Absatzsatz ---
\setlength{\parindent}{0pt}
\setlength{\parskip}{0.5\baselineskip}
\frenchspacing
\clubpenalty=10000
\widowpenalty=10000
\displaywidowpenalty=10000

% U+2212 (echtes Minuszeichen) für pdflatex verfügbar machen
\DeclareUnicodeCharacter{2212}{\textendash}
\DeclareUnicodeCharacter{2192}{\textrightarrow}
% Schwärzung: Volltonbalken in Zeichenbreite
\newcommand{\schwaerzung}{\rule[-0.05em]{0.62em}{0.85em}}
\DeclareUnicodeCharacter{2588}{\schwaerzung}

% Artefakte (### ‹…›) NICHT ins Inhaltsverzeichnis
\setcounter{tocdepth}{2}

% ---------------------------------------------------------------
% Satzbild
% ---------------------------------------------------------------
\usepackage{microtype}
\usepackage{titlesec}

% Akt-Überschriften (#): zentriert, nicht fett, gesperrt
\titleformat{\section}[block]
  {\Large\normalfont\filcenter}{}{0pt}{\textls[180]}
\titlespacing*{\section}{0pt}{1.5\baselineskip}{2.5\baselineskip}

% Kapitel-Überschriften (##): links, fett
\titleformat{\subsection}[block]
  {\large\bfseries}{}{0pt}{}
\titlespacing*{\subsection}{0pt}{1.6\baselineskip}{0.8\baselineskip}

% Artefakt-Kennungen (###)
\titleformat{\subsubsection}[block]
  {\normalsize\bfseries}{}{0pt}{}
\titlespacing*{\subsubsection}{0pt}{0.5\baselineskip}{0.5\baselineskip}

% Inhaltsverzeichnis: Punktführung auch für Akte, gestufte Einzüge,
% großer Titel mit Vorlauf
\makeatletter
\renewcommand*\l@section{\addvspace{0.6\baselineskip}\@dottedtocline{1}{1.5em}{2.3em}}
\renewcommand*\l@subsection{\@dottedtocline{2}{3.8em}{3.2em}}
\renewcommand\tableofcontents{%
  \vspace*{4.5em}%
  {\noindent\Huge\bfseries\contentsname\par}%
  \markboth{\contentsname}{\contentsname}%
  \vspace{2.5em}%
  \@starttoc{toc}%
}
\makeatother
